% Настройки русского языка и кодировок
\usepackage[T2A]{fontenc}
\usepackage[utf8]{inputenc}
\usepackage[russian]{babel}

% Математические пакеты
\usepackage{amsmath, amsfonts, amssymb, amsthm}
\usepackage{geometry}
\usepackage{titlesec}

% Настройка полей (компактные поля для конспекта)
\geometry{left=1cm, right=1cm, top=1.5cm, bottom=1.5cm}

% Настройка заголовков
\titleformat{\section}{\large\bfseries}{}{0em}{}
\titleformat{\subsection}{\normalsize\bfseries}{}{0em}{}

\begin{document}

\begin{center}
    \LARGE \textbf{Методы решения многомерного уравнения теплопроводности}
\end{center}

\section{1. Постановка задачи Коши}
Рассматривается задача Коши для функции $u(x, t) \in C^2(\mathbb{R}^n \times \{t > 0\}) \cap C(\mathbb{R}^n \times \{t \ge 0\})$:
\[
\begin{cases}
u_t = a^2 \Delta u + f(x,t) & (1) \\
u|_{t=0} = u_0(x) & (2)
\end{cases}
\]
\textit{Замечание (Классы единственности):} Решение задачи Коши единственно в классе функций $B_2$, удовлетворяющих оценке $|u(x,t)| \le C e^{\alpha|x|^2}$ для некоторых $C, \alpha > 0$.

\section{2. Общие аналитические формулы (Классика)}
\textbf{Теорема:} Если начальная функция $u_0(x) \in C(\mathbb{R}^n)$ ограничена в $\mathbb{R}^n$, а правая часть $f(x,t) \in C^1(\{t \ge 0\})$ и все ее частные производные до 2-го порядка включительно принадлежат классу ограниченных функций $B_0$, то решение дается \textbf{формулой Пуассона}:
\[ 
u(x,t) = \frac{1}{(2a\sqrt{\pi t})^n} \int\limits_{\mathbb{R}^n} u_0(\xi) e^{-\frac{|\xi-x|^2}{4a^2t}} d\xi + \int\limits_0^t d\tau \int\limits_{\mathbb{R}^n} \frac{f(\xi, \tau)}{(2a\sqrt{\pi(t-\tau)})^n} e^{-\frac{|\xi-x|^2}{4a^2(t-\tau)}} d\xi 
\]

\section{3. Базовая задача с экспонентой и линейное масштабирование}
\textbf{Суть:} Прямое вычисление интеграла Пуассона является громоздким. На практике задачи с экспонентой принято сводить к базовому эталонному случаю.

\vspace{0.2cm}
\noindent\textbf{Базовая (эталонная) задача ($a=1$, показатель $-x^2$):} 
\[
\begin{cases} 
u_t = u_{xx}, \quad x \in \mathbb{R}, \ t>0 \\ 
u|_{t=0} = e^{-x^2}
\end{cases} 
\implies 
\text{Решение:} \quad u(x,t) = \frac{1}{\sqrt{1+4t}} \exp\left(-\frac{x^2}{1+4t}\right)
\]
\textit{Внимание (Экзаменационное требование):} На экзамене этот ответ \textbf{нельзя} писать сразу по памяти. Его \textbf{обязательно} нужно вывести через формулу Пуассона (строгий вывод см. в Приложении).

\vspace{0.2cm}
\noindent\textbf{Линейное масштабирование переменных ($a \neq 1$ или в показателе коэффициент):} \\
Если дана общая задача вида:
\[
\begin{cases} 
u_t = a^2 u_{xx} \\ 
u|_{t=0} = e^{-\alpha^2 x^2}
\end{cases} 
\]
Она сводится к эталонной с помощью замены: \textbf{$\xi = \alpha x, \quad \tau = (a\alpha)^2 t$}. \\
Пусть $u(x,t) = v(\xi, \tau)$. Вычисляя производные по сложной функции ($u_x = \alpha v_\xi$, $u_{xx} = \alpha^2 v_{\xi\xi}$, $u_t = a^2\alpha^2 v_\tau$), исходное уравнение переписывается в виде $v_\tau = v_{\xi\xi}$, а начальное условие — в виде $v|_{\tau=0} = e^{-\xi^2}$. Выполнив обратную замену в эталонном ответе, получается итоговое решение:
\[
u(x,t) = \frac{1}{\sqrt{1+4(a\alpha)^2 t}} \exp\left(-\frac{\alpha^2 x^2}{1+4(a\alpha)^2 t}\right)
\]

\section{4. ПРИЁМ 1: Разбивка на сумму задач}
\textbf{Суть:} В силу линейности дифференциального оператора решение исходной задачи ищется в виде суммы $u = v + w$, где функции отвечают за начальное условие и правую часть по отдельности:

\vspace{0.2cm}
\noindent
\begin{minipage}{0.45\textwidth}
\textbf{A)} $v$ (отвечает за $u_0$): 
\[ \begin{cases} v_t = a^2\Delta v \\ v|_{t=0} = u_0(x) \end{cases} \]
\end{minipage}
\hfill
\begin{minipage}{0.45\textwidth}
\textbf{B)} $w$ (отвечает за $f(x,t)$): 
\[ \begin{cases} w_t = a^2\Delta w + f(x,t) \\ w|_{t=0} = 0 \end{cases} \]
\end{minipage}

\section{5. ПРИЁМ 2: Отщепление гармонического множителя}
\textbf{Суть:} Понижение размерности задачи за счет отделения гармонической части. Применяется, если:
\[
\begin{cases} 
u_t = a^2 \Delta u + g(x_1, \dots, x_m, t) \cdot h(x_{m+1}, \dots, x_n) \\ 
u|_{t=0} = u_0(x_1, \dots, x_m) \cdot h(x_{m+1}, \dots, x_n) 
\end{cases}
\]
Причем $\Delta h = 0$ (функция $h$ является гармонической). \\
\textbf{Решение:} Ищется функция $v(x_1, \dots, x_m, t)$ меньшей размерности из задачи: 
\[
\begin{cases} v_t = a^2 \Delta v + g(x_1, \dots, x_m, t) \\ v|_{t=0} = u_0(x_1, \dots, x_m) \end{cases}
\]
Тогда итоговое решение: $u(x, t) = v(x_1, \dots, x_m, t) \cdot h(x_{m+1}, \dots, x_n)$.

\section{6. ПРИЁМ 3: Использование собственных функций оператора $\Delta$}
\textbf{Скалярный случай:} Если все заданные функции пропорциональны $g(x)$, такой что $\Delta g(x) = \lambda g(x)$: 
\[
\begin{cases} 
u_t = a^2 \Delta u + f(t)g(x) \\ 
u|_{t=0} = \alpha g(x) 
\end{cases} 
\implies \text{Решение ищется в виде } u(x,t) = \varphi(t)g(x).
\]
Функция $\varphi(t)$ находится из решения обыкновенного дифференциального уравнения (ОДУ) первого порядка:
\[ \begin{cases} \varphi'(t) - \lambda a^2 \varphi(t) = f(t) \\ \varphi(0) = \alpha \end{cases} \]

\noindent\textbf{Векторный случай (системы):} Если правые части — комбинации $\{g_1, g_2\}$, где $\begin{pmatrix} \Delta g_1 \\ \Delta g_2 \end{pmatrix} = A \begin{pmatrix} g_1 \\ g_2 \end{pmatrix}$:
\[
\begin{cases}
u_t = a^2 \Delta u + f_1(t)g_1(x) + f_2(t)g_2(x) \\
u|_{t=0} = \alpha_1 g_1(x) + \alpha_2 g_2(x)
\end{cases}
\implies \text{Ищем } u(x,t) = \varphi_1(t)g_1(x) + \varphi_2(t)g_2(x).
\]
Для коэффициентов возникает система ОДУ (где $A^T$ — транспонированная матрица): 
\[
\begin{pmatrix} \varphi_1' \\ \varphi_2' \end{pmatrix} = 
a^2 A^T \begin{pmatrix} \varphi_1 \\ \varphi_2 \end{pmatrix} + \begin{pmatrix} f_1(t) \\ f_2(t) \end{pmatrix}, \qquad
\begin{pmatrix} \varphi_1(0) \\ \varphi_2(0) \end{pmatrix} = \begin{pmatrix} \alpha_1 \\ \alpha_2 \end{pmatrix}
\]

\section{7. ПРИЁМ 4: Полиномиальные данные}
\textbf{Суть:} Если существует такое $N$, что оператор Лапласа в степени $N$ обнуляет начальное условие и правую часть ($\Delta^N u_0 \equiv 0$, $\Delta^N f \equiv 0$), то решение ищется в виде конечного ряда Тейлора по времени. 

\textit{Ограничение:} Правая часть $f$ не должна зависеть от $t$ ($f = f(x)$). \\
Решение задается явной формулой:
\[
u(x,t) = \sum_{k=0}^{N-1} a^{2k} \left[ \frac{t^k}{k!} \Delta^k u_0(x) + \frac{t^{k+1}}{(k+1)!} \Delta^k f(x) \right]
\]

\section{8. ПРИЁМ 5: Сведение к одной переменной (Плоские волны)}
\textbf{Суть:} Если условия зависят от линейной комбинации $\xi = \sum_{i=1}^n \alpha_i x_i$, причем $f(x,t) = \tilde{f}(\xi, t)\big|_{\xi = \sum \alpha_i x_i}$. \\
Решение ищется в виде $u(x,t) = v(\xi, t)$. Уравнение сводится к одномерному уравнению теплопроводности:
\[ 
\begin{cases}
v_t(\xi,t) = \left( a^2 \sum\limits_{i=1}^n \alpha_i^2 \right) v_{\xi\xi}(\xi,t) + \tilde{f}(\xi, t) \\
v|_{t=0} = u_0(\xi)
\end{cases}
\]

\section{9. ПРИЁМ 6: Сферическая симметрия ($n=3$)}
Для пространства $\mathbb{R}^3$, если условия зависят только от радиуса $r = |x|$, решение ищется в виде $u(x,t) = v(r, t)$. Вводится замена \textbf{$w(r,t) = r \cdot v(r,t)$}. Возникает одномерная задача на полуоси:
\[ 
\begin{cases} 
w_t = a^2 w_{rr}, \quad r>0, \ t>0 \\ 
w|_{t=0} = r u_0(r) \\ 
w|_{r=0} = 0 
\end{cases} 
\]
Задача решается методом \textit{нечетного продолжения} начального условия на область $r < 0$. Значение в центре вычисляется как: $u(0,t) = \lim_{r \to 0} \frac{w(r,t)}{r} = w_r(0,t)$.

\section{10. ПРИЁМ 7: Разделение переменных (Произведение решений)}
\textbf{Суть:} Приём \underline{уникален} для уравнения теплопроводности (не применим к волновому!). Используется для однородного уравнения ($f=0$), если начальное условие распадается в произведение:
\[
\begin{cases}
u_t = a^2 \Delta u \\
u|_{t=0} = g(x_1, \dots, x_m) \cdot h(x_{m+1}, \dots, x_n)
\end{cases}
\]
Решение строится как произведение $u(x,t) = v(x_1, \dots, x_m, t) \cdot w(x_{m+1}, \dots, x_n, t)$, где:
\vspace{0.1cm}

\noindent
\begin{minipage}{0.45\textwidth}
\[ \begin{cases} v_t = a^2 \Delta_{x_1 \dots x_m} v \\ v|_{t=0} = g(x_1, \dots, x_m) \end{cases} \]
\end{minipage}
\hfill
\begin{minipage}{0.45\textwidth}
\[ \begin{cases} w_t = a^2 \Delta_{x_{m+1} \dots x_n} w \\ w|_{t=0} = h(x_{m+1}, \dots, x_n) \end{cases} \]
\end{minipage}

\newpage
\appendix
\section{Приложение: Доказательства математических приёмов}

\subsection*{Вывод решения базовой задачи с экспонентой (к Разделу 3)}
\begin{proof}
Запишем формулу Пуассона для эталонной задачи $u_t = u_{xx}$, $u|_{t=0} = e^{-x^2}$ ($a=1$, одномерный случай $n=1$):
\[
u(x,t) = \frac{1}{2\sqrt{\pi t}} \int\limits_{-\infty}^{+\infty} e^{-\xi^2} e^{-\frac{(\xi-x)^2}{4t}} d\xi = \frac{1}{2\sqrt{\pi t}} \int\limits_{-\infty}^{+\infty} \exp\left( -\xi^2 - \frac{(\xi-x)^2}{4t} \right) d\xi.
\]
Преобразуем показатель степени, раскрыв скобки и сгруппировав слагаемые относительно $\xi$:
\[
-\xi^2 - \frac{\xi^2 - 2\xi x + x^2}{4t} = -\frac{4t\xi^2 + \xi^2 - 2\xi x + x^2}{4t} = -\frac{(1+4t)\xi^2 - 2\xi x + x^2}{4t}.
\]
Выделим полный квадрат по переменной интегрирования $\xi$:
\[
-\frac{1+4t}{4t} \left( \xi^2 - 2\xi \frac{x}{1+4t} \right) - \frac{x^2}{4t} = -\frac{1+4t}{4t} \left( \xi - \frac{x}{1+4t} \right)^2 + \frac{1+4t}{4t} \frac{x^2}{(1+4t)^2} - \frac{x^2}{4t}.
\]
Свободный от $\xi$ член (константа по отношению к интегралу) упрощается:
\[
\frac{x^2}{4t(1+4t)} - \frac{x^2(1+4t)}{4t(1+4t)} = \frac{x^2 - x^2 - 4tx^2}{4t(1+4t)} = -\frac{x^2}{1+4t}.
\]
Таким образом, интеграл принимает вид:
\[
u(x,t) = \frac{1}{2\sqrt{\pi t}} e^{-\frac{x^2}{1+4t}} \int\limits_{-\infty}^{+\infty} \exp\left( -\frac{1+4t}{4t} \left( \xi - \frac{x}{1+4t} \right)^2 \right) d\xi.
\]
Сделаем линейную замену переменной $\eta = \sqrt{\frac{1+4t}{4t}} \left( \xi - \frac{x}{1+4t} \right)$, откуда $d\xi = \sqrt{\frac{4t}{1+4t}} d\eta = \frac{2\sqrt{t}}{\sqrt{1+4t}} d\eta$. Пределы интегрирования при этом остаются бесконечными:
\[
u(x,t) = \frac{1}{2\sqrt{\pi t}} e^{-\frac{x^2}{1+4t}} \frac{2\sqrt{t}}{\sqrt{1+4t}} \int\limits_{-\infty}^{+\infty} e^{-\eta^2} d\eta.
\]
Используя табличный интеграл Эйлера-Пуассона $\int_{-\infty}^{+\infty} e^{-\eta^2} d\eta = \sqrt{\pi}$, окончательно получаем:
\[
u(x,t) = \frac{1}{\sqrt{\pi}} e^{-\frac{x^2}{1+4t}} \frac{1}{\sqrt{1+4t}} \cdot \sqrt{\pi} = \frac{1}{\sqrt{1+4t}} \exp\left(-\frac{x^2}{1+4t}\right).
\]
\end{proof}

\subsection*{Доказательство к Приёму 1 (Разбивка на сумму задач)}
\begin{proof}
Пусть функции $v, w$ являются решениями соответствующих задач A и B. Требуется доказать, что их сумма $u(x,t) = v(x,t) + w(x,t)$ является решением исходной задачи. В силу линейности оператора дифференцирования:
\[
u_t = v_t + w_t = a^2\Delta v + (a^2\Delta w + f(x,t)) = a^2\Delta(v+w) + f(x,t) = a^2\Delta u + f(x,t).
\]
Начальные условия: $u|_{t=0} = v|_{t=0} + w|_{t=0} = u_0(x) + 0 = u_0(x)$. По теореме единственности, найденное решение является искомым.
\end{proof}

\subsection*{Доказательство к Приёму 2 (Отщепление гармонического множителя)}
\begin{proof}
Оператор Лапласа от произведения раскрывается как $\Delta(v \cdot h) = (\Delta v)h + 2(\nabla v, \nabla h) + v(\Delta h)$.
Так как $v$ и $h$ зависят от разных переменных, $(\nabla v, \nabla h) = 0$. Так как $h$ гармоническая, $\Delta h = 0$. Отсюда $\Delta(v \cdot h) = (\Delta v) \cdot h$.
Домножая уравнение для $v$ на $h$, получаем:
\[
v_t \cdot h = (a^2\Delta v + g) \cdot h \implies (v \cdot h)_t = a^2(\Delta v) \cdot h + g \cdot h \implies u_t = a^2\Delta u + g \cdot h.
\]
\end{proof}

\subsection*{Доказательство к Приёму 3* (Собственные функции оператора $\Delta$)}
\begin{proof}
Ищем решение в виде $u(x,t) = \varphi_1(t)g_1(x) + \varphi_2(t)g_2(x)$. По условию задана матрица $A = \begin{pmatrix} a_{11} & a_{12} \\ a_{21} & a_{22} \end{pmatrix}$, такая что $\Delta g_1 = a_{11} g_1 + a_{12} g_2$ и $\Delta g_2 = a_{21} g_1 + a_{22} g_2$. Подставляем $u$ в уравнение $u_t = a^2 \Delta u + f_1 g_1 + f_2 g_2$:
\[
\varphi_1' g_1 + \varphi_2' g_2 = a^2 \left[ \varphi_1(a_{11} g_1 + a_{12} g_2) + \varphi_2(a_{21} g_1 + a_{22} g_2) \right] + f_1 g_1 + f_2 g_2.
\]
Перегруппируем слагаемые при базисных функциях $g_1$ и $g_2$:
\[
\begin{cases}
\varphi_1' = a^2(a_{11}\varphi_1 + a_{21}\varphi_2) + f_1 \\
\varphi_2' = a^2(a_{12}\varphi_1 + a_{22}\varphi_2) + f_2
\end{cases}
\]
В правой части сформировалась матрица $A^T$. В матричной форме это записывается как $\vec{\varphi}' = a^2 A^T \vec{\varphi} + \vec{f}$.
\end{proof}

\subsection*{Доказательство к Приёму 4 (Полиномиальные данные)}
\begin{proof}
Рассмотрим случай с начальным отклонением $u_0(x)$ при $f \equiv 0$ (доказательство для $f(x)$ аналогично). Подставим ряд $u = \sum_{k=0}^{N-1} a^{2k} \frac{t^k}{k!} \Delta^k u_0(x)$ в уравнение $u_t = a^2\Delta u$.
Вычислим производную по $t$ (при $k=0$ производная константы равна 0):
\[
u_t = \sum_{k=1}^{N-1} a^{2k} \frac{k \cdot t^{k-1}}{k!} \Delta^k u_0(x) = \sum_{k=1}^{N-1} a^{2k} \frac{t^{k-1}}{(k-1)!} \Delta^k u_0(x).
\]
Сделаем сдвиг индекса $l = k - 1$:
\[
u_t = \sum_{l=0}^{N-2} a^{2l+2} \frac{t^l}{l!} \Delta^{l+1} u_0(x).
\]
Вычислим $a^2\Delta u$:
\[
a^2\Delta u = a^2 \sum_{k=0}^{N-1} a^{2k} \frac{t^k}{k!} \Delta^{k+1} u_0(x) = \sum_{k=0}^{N-1} a^{2k+2} \frac{t^k}{k!} \Delta^{k+1} u_0(x).
\]
Оба ряда совпадают для индексов от $0$ до $N-2$. В ряде для $a^2\Delta u$ остается слагаемое при $k = N-1$:
\[ a^{2N} \frac{t^{N-1}}{(N-1)!} \Delta^N u_0(x). \]
По условию $\Delta^N u_0 \equiv 0$, поэтому это слагаемое равно нулю. Тождество $u_t = a^2\Delta u$ доказано.
\end{proof}

\subsection*{Доказательство к Приёму 5 (Плоские волны)}
\begin{proof}
Пусть $u(x,t) = v(\xi, t)$, где $\xi = \sum_{i=1}^n \alpha_i x_i$. По правилу дифференцирования сложной функции:
$u_t = v_t$, $u_{x_i} = v_\xi \cdot \alpha_i$, $u_{x_i x_i} = v_{\xi\xi} \cdot \alpha_i^2$. Оператор Лапласа:
\[
\Delta u = \sum_{i=1}^n u_{x_i x_i} = \sum_{i=1}^n \left( v_{\xi\xi} \cdot \alpha_i^2 \right) = v_{\xi\xi} \sum_{i=1}^n \alpha_i^2.
\]
Подставляя в уравнение $u_t = a^2\Delta u + \tilde{f}$, получаем: $v_t = a^2 \left(\sum_{i=1}^n \alpha_i^2 \right) v_{\xi\xi} + \tilde{f}(\xi, t)$.
\end{proof}

\subsection*{Доказательство к Приёму 6 (Сферическая симметрия)}
\begin{proof}
Для сферически симметричной функции $\Delta v = v_{rr} + \frac{2}{r}v_r$. Подстановка $w(r,t) = r \cdot v(r,t) \implies v = \frac{w}{r}$.
\[
v_t = \frac{w_t}{r}, \quad v_r = \frac{w_r}{r} - \frac{w}{r^2}, \quad v_{rr} = \frac{w_{rr}}{r} - \frac{2w_r}{r^2} + \frac{2w}{r^3}.
\]
Подставляем в выражение для Лапласиана:
\[
\Delta v = \left( \frac{w_{rr}}{r} - \frac{2w_r}{r^2} + \frac{2w}{r^3} \right) + \frac{2}{r} \left( \frac{w_r}{r} - \frac{w}{r^2} \right) = \frac{w_{rr}}{r}.
\]
Из уравнения $v_t = a^2\Delta v$ получаем $\frac{w_t}{r} = a^2 \frac{w_{rr}}{r} \implies w_t = a^2 w_{rr}$. Для ограниченности $v$ при $r \to 0$ требуется $w(0, t) = 0$.
\end{proof}

\subsection*{Доказательство к Приёму 7 (Разделение переменных)}
\begin{proof}
Рассмотрим функцию $u = v \cdot w$. Вычислим ее производную по времени:
\[ u_t = (v \cdot w)_t = v_t \cdot w + v \cdot w_t. \]
Подставим уравнения $v_t = a^2 \Delta v$ и $w_t = a^2 \Delta w$:
\[ u_t = (a^2 \Delta v) \cdot w + v \cdot (a^2 \Delta w) = a^2 (\Delta v \cdot w + v \cdot \Delta w). \]
Так как функции $v$ и $w$ зависят от независимых наборов пространственных переменных, их градиенты ортогональны $(\nabla v, \nabla w) = 0$. Следовательно, оператор Лапласа от произведения раскрывается точно так же: $\Delta(v \cdot w) = \Delta v \cdot w + v \cdot \Delta w$.
Отсюда $u_t = a^2 \Delta(v \cdot w) = a^2 \Delta u$. Начальное условие $u|_{t=0} = v|_{t=0} \cdot w|_{t=0} = g \cdot h$ выполняется очевидным образом.
\end{proof}

\end{document}